\section{Contract2Video}
\label{sec:method}

Contract2Video separates language-level planning from evidence about rendered
media. Its three stages are contract compilation, specialist planning with
abstaining arbitration, and rendering with dependency-aware repair. Planning
agents propose what a video should contain. Observers assess what its pixels
actually show. The latter assessment alone can justify output acceptance.
Figure~\ref{fig:architecture} will illustrate this distinction.

\subsection{Task and Acceptance Semantics}
\label{sec:task}

Given a request $x$, the system produces an ordered video
$V=(v_1,\ldots,v_m)$ under a generation-call budget $B_N$ and a cost budget
$B_C$. A requirement is a predicate $c_j$ whose scope
$\Omega_j\subseteq\{1,\ldots,m\}$ identifies the relevant shots.
The contract partitions requirements into hard constraints $\mathcal H$ and
soft preferences $\mathcal S$. We distinguish the latent semantic condition
\begin{equation}
 H(V;\mathcal C)=\prod_{j\in\mathcal H}c_j(V_{\Omega_j})
 \label{eq:true-acceptance}
\end{equation}
from its operational approximation $\widehat H(V;\mathcal C)$, computed by
fallible observers. The input contract alone cannot satisfy
Equation~\ref{eq:true-acceptance}. The arguments must include rendered media.
Aesthetic quality is evaluated separately and does not override an unsupported
hard requirement. The primary setting is silent multi-shot video. Audio
relations remain representable but are inactive in the primary experiments.

The intended optimization is to increase hard-constraint satisfaction and
perceptual quality within fixed resource limits. We do not assume access to
the true predicates during generation, nor claim to solve a globally optimal
video-editing problem. Our controller implements a conservative, measurable
decision procedure whose semantic accuracy must be established empirically.

\subsection{Stage I: Contract Compilation}
\label{sec:contract}

Following the architecture, the semantic contract is
\begin{equation}
 \mathcal C=(E,A,O,R_T,R_S,R_A,\mathcal H,\mathcal S),
\end{equation}
where $E$ denotes persistent characters or other named entities, $A$ their
attributes, $O$ scene objects, and $R_T$, $R_S$, and $R_A$ temporal, spatial,
and audio relations. These components are serialized as stable entity
identifiers and typed constraints. Objects and characters share the same
identifier namespace. An object's role is expressed in its description and
constraints. In particular, the JSON interface does not require an object
detector to recognize every entity type before planning begins.

The compiler also emits a shot plan
$\mathcal P=(Q,D,\Sigma)$, with ordered shots
$Q=(q_1,\ldots,q_m)$, a dependency DAG $D$, and planned boundary states
$\Sigma$. Each shot specifies a prompt, duration, participating entities,
and input/output states. An edge $i\!\rightarrow\!k$ means that rendering
shot $k$ may depend on shot $i$ through a reference asset or a propagated
state. Shot order is a topological order. Narrative events are represented
through action constraints and state transitions attached to these shots.
we do not equate a textual state with an observed event.

Each serialized constraint stores its identifier, kind, natural-language
description, shot scope, hardness, and positive evaluation weight. Supported
kinds are entity, attribute, action, spatial, temporal, identity, text, and
audio. A structural validator checks identifier uniqueness, reference
resolution, nonempty scopes, positive finite durations and weights, and
backward-only dependencies. These checks establish structural validity,
not consistency with the user's intent. Compiler fidelity is therefore
evaluated separately against human-annotated requirements.

The contract is frozen before comparing repair policies for a given request.
This controls for differences introduced by prompt decomposition. A separate
end-to-end condition includes compiler errors. Neither repair nor arbitration
may delete a hard constraint, lower its hardness, or replace it with an easier
requirement merely to obtain acceptance.

\subsection{Stage II: Specialist Planning and Arbitration}
\label{sec:specialists}

Three specialists cover complementary planning responsibilities. The story
planner proposes event order and shot boundaries. The identity specialist
tracks persistent descriptions, object ownership, and cross-shot references.
The camera specialist proposes framing and motion while preserving the
visibility needed to check requirements. An audio specialist is reserved for
the audiovisual extension. Specialists receive the frozen contract and
propose patches to the shot plan, rather than independently rewriting the
user's requirements.

Before rendering, an arbiter can reject a plan with a demonstrable symbolic
conflict, such as a dependency cycle or incompatible declared object states.
It can also request clarification of an uncheckable constraint. However,
the absence of visual evidence before rendering is expected. It does not
justify a visual violation verdict. Pre-render approval means only that a
plan passes the available symbolic checks. Specialist disagreement after
rendering is resolved by media evidence as described below, not by the
number of agents repeating a claim. Multiple calls to one model are treated
as potentially correlated observations.

\subsection{Timestamped Evidence and Three-Way Decisions}
\label{sec:evidence}

For constraint $j$, an observer emits records
\begin{equation}
 e=(j,z,r,\Omega_e,T_e,h_e,u_e),
\end{equation}
where $z\in\{+, -, ?\}$ denotes support, contradiction, or abstention.
The value $r\in[0,1]$ is an observation score. $\Omega_e$ lists the referenced
shots, while $T_e$ contains within-shot time intervals. The hash $h_e$ binds
those shots to their content hashes, and $u_e$ identifies the observation
source. A rationale is
stored for inspection but is not itself visual evidence. The score is not
assumed to be a calibrated probability of truth.

Every non-abstaining record must name an existing constraint, reference only
shots in its scope, and provide a valid interval and matching content hash
for every referenced asset. Invalid records become unknown. Hash validation
prevents evidence for an older render from silently being assigned to a
replacement. It does not prove that an observer's interpretation is correct.
Likewise, valid timestamps establish localization metadata, not perception
accuracy.

Let $\mathcal E_j^+$ and $\mathcal E_j^-$ be valid support and contradiction
records whose scores exceed their direction-specific thresholds. Let
$b_j=\mathbf 1[\Omega_j\subseteq\bigcup_{e\in\mathcal E_j^+}\Omega_e]$
indicate complete support coverage. The conservative aggregation rule is
\begin{equation}
 d_j =
 \begin{cases}
 \textsc{violated}, & \mathcal E_j^-\neq\varnothing,\\
 \textsc{supported}, & \mathcal E_j^-=\varnothing,\ b_j=1,\\
 \textsc{unknown}, & \text{otherwise}.
 \end{cases}
 \label{eq:arbitration}
\end{equation}
Thus an accepted contradiction takes precedence over support, and incomplete
scope coverage cannot pass. For an identity or temporal relation, the
observer must assess the relation jointly across its scope. Observing its
individual endpoints is insufficient. Scope coverage is a necessary
bookkeeping condition, not a substitute for that semantic check. Joint
relation assessment is a required integration and evaluation criterion for
the observer interface.

Let $a(V)$ indicate that every required asset exists and is decodable.
The operational acceptance condition is
\begin{equation}
 \widehat H(V;\mathcal C)
 =a(V)
  \prod_{j\in\mathcal H}\mathbf 1[d_j=\textsc{supported}].
 \label{eq:operational-acceptance}
\end{equation}
An accepted output is therefore \emph{supported by the configured verifier},
not guaranteed correct. Missing evidence yields abstention. It blocks
acceptance of a hard constraint, but is reported separately from an observed
violation. This distinction matters when visibility is poor or an observer
cannot judge a subtle motion.

\subsection{Threshold Selection and Certification}
\label{sec:calibration}

We select thresholds using a development split and assess them on a separate
certification split with disjoint case identifiers. A calibration label
indicates whether an observer's directional verdict is correct, rather than
whether the video is globally satisfactory. For each predeclared stratum
$g$ (support or contradiction in the primary setting), threshold $\tau_g$
is selected to retain correct verdicts subject to a development false-positive
constraint. It is frozen before certification.

Let $n_g$ denote the number of incorrect candidate verdicts in the
certification split, and let $k_g$ of these exceed $\tau_g$. The negative
acceptance rate is estimated by $k_g/n_g$. For $K$ fixed strata, the
implemented upper bound is
\begin{equation}
 U_g=\min\left\{1,\frac{k_g}{n_g}
       +\sqrt{\frac{\log(K/\alpha)}{2n_g}}\right\}.
 \label{eq:hoeffding}
\end{equation}
With independent certification examples within each stratum, fixed
thresholds, and the same sampling distribution, Hoeffding's inequality and
a union bound give simultaneous coverage at least $1-\alpha$. A stratum
with no negative examples is uncertified. The strict policy abstains if
required certification is unavailable or if any required bound exceeds its
target $\epsilon$.

This is a bound on $\Pr(r\geq\tau_g\mid\text{verdict incorrect})$.
It is not a bound on the error fraction among accepted verdicts, the error
probability of a complete video, or the risk after adaptive repair.
Repeated observations from one case must not be counted as independent
certification examples. In the proposed setting, one preselected
candidate per case and direction is used for certification. Richer repeated
observations remain available for descriptive analysis. Appendix
Table~\ref{tab:calibration} recomputes the bound from the retained sample-size
and error summaries. These entries do not certify a 5\% error target.
Integer error counts, eligible negative-example counts, and case identifiers
are required before any empirical certification claim.

\subsection{Stage III: Dependency-Aware Local Repair}
\label{sec:repair}

After a render wave, all constraints are reassessed against the current
assets. Let $\mathcal U_t$ be the hard constraints marked violated or unknown
at round $t$. Their shot scopes define repair roots
\begin{equation}
 R_t=\bigcup_{j\in\mathcal U_t}\Omega_j.
\end{equation}
The repair set contains these roots and their transitive dependents:
\begin{equation}
 A_t=R_t\cup\operatorname{Desc}_{D}(R_t).
 \label{eq:repair-closure}
\end{equation}
Shots in $A_t$ are regenerated in timeline order. All other assets are
retained byte-for-byte. Descendant repair prevents a later shot from retaining
an obsolete reference to a changed predecessor. Closure can equal the entire
timeline, so the method does not promise a saving on every request.

For observed violations, repair feedback identifies the constraint, affected
span, and counterexample description. For unknown decisions, feedback instead
states which evidence is unavailable. It must not invent a counterexample.
A repair adapter may revise the shot prompt or regenerate a reference while
preserving the frozen contract. Prompt changes and reference consumption are
logged explicitly. Re-rendering with a new seed alone is a distinct baseline.
The rule in Equation~\ref{eq:repair-closure} is a deterministic heuristic,
not an exact minimum-cost edit operator and not a guarantee that the next
attempt will succeed.

Each attempted render, including a failed call, consumes the call budget.
Before a request, the controller checks whether its conservative cost
estimate fits the remaining budget. A monetary hard cap holds only if
pre-call estimates upper-bound billed costs. Unknown provider charges cannot
be bounded by bookkeeping alone. The system terminates when all hard
constraints are supported, a budget or round limit prevents further work,
or an execution error occurs. It records accepted, unresolved,
budget-exhausted, and failed outcomes separately. A composed video may be
retained for an unresolved run, with its status preserved.

Table~\ref{tab:controller-pseudocode} makes the execution order explicit.
The controller renders versioned assets, re-observes the current timeline after
each wave, and computes the transitive dependency closure before repair. The
contract remains fixed throughout this process.

\begin{table}[t]
\centering
\tablefont
\setlength{\tabcolsep}{4pt}
\renewcommand{\arraystretch}{1.08}
\begin{tabularx}{\columnwidth}{@{}>{\raggedleft\arraybackslash}p{4.0em}>{\raggedright\arraybackslash}X@{}}
\toprule
\rowgreendeep
\multicolumn{2}{@{}l@{}}{\textbf{Contract2Video controller}}\\
\rowgreenlight
\textbf{Input} & Request $x$, budgets $(B_N,B_C)$, thresholds $\tau$, renderer $\mathsf{G}$, and observer $\mathsf{O}$.\\
\rowgreenlight
\textbf{Output} & Video $V$, status $s$, and execution trace $T$.\\
\midrule
\rowgreenmid
\multicolumn{2}{@{}l@{}}{\textbf{1. Compile and validate}}\\
\rowgreenpale
1 & $(\mathcal C,\mathcal P)\leftarrow\mathsf{Compile}(x)$ with $\mathcal P=(Q,D,\Sigma)$.\\
\rowgreenpale
2 & \textbf{if} $\neg\mathsf{Valid}(\mathcal C,\mathcal P)$ \textbf{return} $(\varnothing,\textsc{failed},T)$.\\
\rowgreenpale
3 & $W\leftarrow Q$ in topological order, $n\leftarrow0$, and $b\leftarrow0$. Initialize versioned assets and $T$.\\
\midrule
\rowgreenmid
\multicolumn{2}{@{}l@{}}{\textbf{2. Render, observe, and decide}}\\
\rowgreenpale
4 & \textbf{while} $W\neq\varnothing$ and $\mathsf{BudgetOK}(n,b,B_N,B_C)$ \textbf{do}\\
\rowgreenpale
5 & \quad \textbf{for} $q\in W$ in topological order \textbf{do}\\
\rowgreenpale
6 & \qquad render $v_q\leftarrow\mathsf{G}(q,\mathcal C,\Sigma,T)$, increment $n,b$ before the call, and store its version and content hash.\\
\rowgreenpale
7 & \quad $\mathcal E\leftarrow\mathsf{O}(V_{\mathrm{current}})$. Invalidate records with an invalid scope, interval, or hash.\\
\rowgreenpale
8 & \quad $d_j\leftarrow\mathsf{Decide}(\mathcal E_j,\tau_j)$ for each $j\in\mathcal H$ using contradiction-first aggregation.\\
\rowgreenpale
9 & \quad \textbf{if} $a(V_{\mathrm{current}})=1$ and all $d_j=\textsc{supported}$ \textbf{return} $(V_{\mathrm{current}},\textsc{accepted},T)$.\\
\rowgreenpale
10 & \quad $U\leftarrow\{j\in\mathcal H:d_j\in\{\textsc{violated},\textsc{unknown}\}\}$.\\
\rowgreenpale
11 & \quad $R\leftarrow\bigcup_{j\in U}\Omega_j$ and $A\leftarrow R\cup\operatorname{Desc}_{D}(R)$.\\
\rowgreenpale
12 & \quad $F\leftarrow\mathsf{Explain}(U,\mathcal E)$. Revise prompts or references only for $A$, preserving $\mathcal C$.\\
\rowgreenpale
13 & \quad $W\leftarrow A$ in topological order and append the current-wave records to $T$.\\
\rowgreenpale
14 & \textbf{end while}\\
\midrule
\rowgreenlight
15 & \textbf{if} $\neg\mathsf{BudgetOK}(n,b,B_N,B_C)$ \textbf{return} $(V_{\mathrm{current}},\textsc{budget\mbox{-}exhausted},T)$. Otherwise return $(V_{\mathrm{current}},\textsc{unresolved},T)$.\\
\bottomrule
\end{tabularx}
\caption{Controller pseudocode for contract compilation, evidence-grounded acceptance, and dependency-aware repair. Every render attempt is charged before dispatch. Invalid evidence yields \textsc{unknown}, and repair preserves the frozen contract.}
\label{tab:controller-pseudocode}
\end{table}

\paragraph{Implementation boundary.}
The available controller implements schema validation, evidence validation
and aggregation, threshold selection/certification, dependency closure,
budgets, versioned assets, and trace export. The generic command adapter
exposes feedback and reference arguments. The current direct Wan adapter
performs prompt-based generation and does not yet consume those arguments.
Its repeated calls correspond to resampling. Provider-backed visual
observation, full specialist planning, and counterexample-conditioned
rendering remain integration tasks for the proposed experiments. The paper
does not treat deterministic examples or interface tests as evidence of
video quality.
